% Corpus document: a non-Latin-script paper (DESIGN.md §8). fontspec and polyglossia over
% XeTeX, with Greek, Cyrillic and Hebrew (right to left) in one document: font loading and
% script shaping are their own cost, and a font the bundle lacks is its own failure class.
\documentclass[11pt]{article}
\usepackage[margin=2.5cm]{geometry}
\usepackage{fontspec}
\usepackage{polyglossia}
\setdefaultlanguage{english}
\setotherlanguages{greek,russian,hebrew}
% Fonts in Tectonic's bundle are found by file name, not family name: `DejaVu Serif` is not
% found, `FreeSerif.otf` is (fixtures/corpus/README.md). GNU FreeFont covers all three scripts.
\setmainfont{FreeSerif.otf}[BoldFont=FreeSerifBold.otf, ItalicFont=FreeSerifItalic.otf]
\setsansfont{FreeSans.otf}
\newfontfamily\greekfont{FreeSerif.otf}
\newfontfamily\cyrillicfont{FreeSerif.otf}
\newfontfamily\hebrewfont{FreeSerif.otf}[Script=Hebrew]

\begin{document}
\section{Three scripts, one document}

An English paragraph first, so the default language is exercised before the others.

\begin{greek}
Η τοπική μεταγλώττιση κάνει τον συντάκτη γρήγορο: κάθε αλλαγή φαίνεται στο PDF σε λιγότερο
από ένα δευτερόλεπτο, χωρίς ουρά αναμονής και χωρίς διακομιστή.
\end{greek}

\begin{russian}
Локальная компиляция делает редактор быстрым: каждое изменение появляется в PDF меньше чем
за секунду, без очереди и без сервера.
\end{russian}

\begin{hebrew}
הידור מקומי הופך את העורך למהיר: כל שינוי מופיע בקובץ תוך פחות משנייה, בלי תור ובלי שרת.
\end{hebrew}

\section{Mixed within a line}
The word \textgreek{ταχύτητα} (speed), the word \textrussian{скорость}, and the word
\texthebrew{מהירות} all mean the same thing.

\end{document}
